\documentclass{stex}

\libusepackage{naproche}
\libinput{preamble}

\begin{document}
\begin{smodule}{limit-ordinals.ftl}

\importmodule[libraries/set-theory]{zero.ftl}
\importmodule[libraries/set-theory]{successor-ordinals.ftl}

\symdef{Limit}{\mathop{\textsf{lim}}}

\symdecl*{limit ordinal}

\begin{convention*}[forthel,for=Limit]
  Let $\Limit x$ stand for $\UnionOver x$.
\end{convention*}

\begin{proposition*}[forthel]
  Let $x$ be a subset of $\Ordinals$.
  Then $\Limit x$ is an ordinal.
\end{proposition*}
\begin{proof}[forthel]
  (1) $\Limit x$ is transitive.
  \begin{proof}
    Let $y \Elem \Limit x$ and $z \Elem y$.
    Take $w \Elem x$ such that $y \Elem w$.
    Hence $w$ is transitive.
    Thus $z \Elem w$.
    Therefore $z \Elem \Limit x$.
  \end{proof}

  (2) Every element of $\Limit x$ is transitive.
  \begin{proof}
    Let $y \Elem \Limit x$.
    Let $z \Elem y$ and $v \Elem z$.
    Take $w \Elem x$ such that $y \Elem w$.
    Hence $w$ is an ordinal.
    Thus $y$ is an ordinal.
    Therefore $y$ is transitive.
    Consequently $v \Elem y$.
  \end{proof}
\end{proof}

\begin{definition*}[forthel,for=limit ordinal]
  A \emph{limit ordinal} is an ordinal $\lambda$ such that neither $\lambda$ is a successor ordinal nor $\lambda \Eq \Zero$.
\end{definition*}

\end{smodule}
\end{document}
